\chapter{相关技术与问题分析}\label{chap:related}

\section{ART 执行引擎的执行链路与关键组件}

ART作为完整的执行系统，涵盖应用装载、DEX 校验、解释执行、即时编译、安装期编译、JNI 桥接和运行时辅助服务等多个核心模块\citep{AndroidRuntime2024, ArtJitCompiler2024}。从源码组织架构分析，\texttt{runtime/} 负责管理线程、对象模型、类链接、垃圾回收及运行时服务，\texttt{interpreter/} 负责基础解释执行路径，\texttt{compiler/}、\texttt{jit/} 与 \texttt{dex2oat/} 共同构建多级编译基础设施，\texttt{libdexfile/} 承担 DEX 文件的解析与校验职责。针对新架构平台的适配，需将上述组件视作闭合的执行链进行整体移植与协同验证。

基于应用执行流程，目标方法在通过类加载与 DEX 校验后，首先进入 C++ 解释器或 Nterp 汇编解释器等冷路径执行；待方法热度累积至预设阈值，ART 将动态触发 baseline JIT、优化重编译或栈上替换（OSR）等高性能执行路径\citep{ArtJitCompiler2024, Fink2003, Lameed2013}。此外，字符串处理、对象分配、字段访问、异常抛出及 JNI 调用等高频操作，高度依赖运行时快速入口点（Quick Entrypoints）与内联辅助机制（Intrinsics）。因此，平台适配的核心在于使解释器、编译器后端、运行时入口点与标准库机制在目标架构上形成一致的控制流、数据流与异常处理语义。

上述执行链构成了本文后续研究的系统边界。第 3 章聚焦于基础执行能力的链路构建，重点验证解释器、JIT、JNI 与运行时等核心组件的连通性与正确性；第 5 章与第 6 章则在建立有效执行环境的基础上，深入分析并优化特定的性能瓶颈。本章将重点剖析与后续架构适配及性能调优强相关的执行机制与架构特征。

\section{LoongArch 平台约束与问题映射}

LoongArch 体系结构特征直接影响目标平台的运行时性能表现。结合本文的研究问题，本节重点提炼三类与运行时机制强相关的架构约束：控制流与 JIT 编译约束、数据传输与访存约束，以及 LSX 向量化与后端代码映射（Lowering）约束\citep{LoongArchISA2023, LoongArchISAv2, LoongArchELFABI2025}。

\begin{figure}[htbp]
    \centering
    \includegraphics[width=1.00\textwidth]{related/constraint-problem-mapping.pdf}
    \caption{LoongArch 平台关键约束、主导成本与后续优化问题的映射关系}
    \label{fig:constraint-problem-mapping}
\end{figure}

如图~\ref{fig:constraint-problem-mapping} 所示，本文将架构约束与性能调优问题建立对应映射：控制流约束对应 JIT 触发时机的判定，数据访存约束涉及字符串热路径的连续内存传输，而 LSX 向量化特性则决定了单指令多数据流（SIMD）编译链路的生成机制。该映射关系为后续探索场景依赖的性能优化提供了底层架构支撑。

\subsection{控制流与 JIT 编译相关的架构约束}

LoongArch 采用定长 32 位指令格式，基于严格的寄存器调用约定、分支编码与地址构造规则组织控制流\citep{LoongArchISA2023, LoongArchELFABI2025}。在 ART 执行上下文中，上述机制不仅涉及常规 ABI 兼容问题，更直接影响解释器分派、慢路径跳转、跳板函数（Trampoline）入口配置、异常栈恢复及 JIT 生成代码的回填逻辑。若寄存器传参、返回值通道、长短跳转或 PC 相对地址构造处理不当，将导致整个执行链路的语义破坏。

此类架构约束直接影响动态编译策略的制定。短生命周期工作负载对解释器驻留时间、JIT 启动时机及编译固定成本高度敏感；且早期编译的性能收益取决于目标平台在构建入口、修补链接、维护调试信息与异常边界时的底层开销。因此，控制流操作的架构代价与工作负载生命周期的耦合程度，构成了即时编译启发式策略的核心限制条件。

\subsection{数据搬移与访存相关的架构约束}

LoongArch 的运行时执行流程受制于显式访存（Load/Store）指令吞吐、地址计算延时及边界检查机制\citep{LoongArchISA2023, LoongArchISAv2}。在运行时热路径中，数据拷贝的开销不仅取决于传输数据量，亦受制于循环控制、边界检查、地址更新、对象布局解析及安全语义维护。例如，在字符串构造、复制与编码转换等高频操作中，上述微指令级开销易在密集型负载下累积为显著的性能瓶颈\citep{Haubl2008OptimizedStrings, Haubl2010CompactStrings}。

针对此类访存约束，若运行时采用细粒度的逐元素搬运机制，将引入密集的地址计算与控制转移开销。通过在保障对象语义一致性的前提下采用内存批量传输，可有效降低热路径中的数据流局部开销，进而提升整体吞吐量。

\subsection{LSX 向量能力、对齐约束与后端差异}

LoongArch64 提供 LSX 128 位向量扩展，配备独立的向量寄存器及数据并行计算指令支持\citep{LoongArchISA2023, LoongArchISAv2}。然而，硬件指令支持需辅以完备的动态编译器后端实现方可生效。在混合编译框架中，向量化转换需经历依赖分析、类型检查、寄存器分配及目标指令映射（Lowering）等环节\citep{Nie2010VectorizationJava}。任何一环的验证缺失均会导致代码生成回退至标量执行路径。

鉴于后端工程实现的复杂性，针对规则浮点数值计算，构建轻量且闭环的显式代码生成链路具有现实意义。本研究将探讨 ART 向量中间表示（IR）向 LoongArch 汇编后端映射的核心链路构建条件，并结合微架构特征界定该优化的生效边界与回退机制。

\section{JIT、运行时与向量化的原理基础}

动态编译优化的核心在于通过运行期的编译开销换取更高的机器码质量\citep{Arnold2005, Paleczny2001, Suganuma2001}。JIT 触发时机的设定需综合权衡方法生命周期、解释执行时间、编译开销及热点演进速率。对于长生命周期负载，早期编译可加速摊销固定成本；而针对短生命周期负载，激进的编译策略则可能导致系统资源竞争与性能倒挂。该机制为探索适配特定生命周期的动态编译参数调优提供了理论依据。

运行时高频路径优化则侧重于缩短高频操作的底层指令路径，包括消除细粒度控制流、简化辅助函数调用栈，以及引入内联函数（Intrinsic）与快速路径实现\citep{Haubl2008OptimizedStrings, Haubl2010CompactStrings, ArtJitCompiler2024}。此类优化方案在特定执行特征被稳定命中时，能够针对特定高频操作提供确定性且边界清晰的性能提升。

向量化优化的前提是目标循环具备规则的迭代结构、清晰的数据依赖及兼容的通道（Lane）语义，从而允许编译器将标量 IR 提升为向量 IR 并映射至底层 SIMD 指令\citep{Nie2010VectorizationJava}。因此，安全映射 IR 节点至 LSX 指令、处理数据对齐约束及界定标量回退条件，构成了数据并行优化落地的理论基础。

\section{工作负载特征与主导成本分析}

基准测试工作负载的性能瓶颈需结合系统底层开销类型进行评估\citep{Blackburn2006, SPECjvm2008, Lengauer2017GCBench, Prokopec2019Renaissance}。本文选用多组具备典型特征的测试集，以隔离并度量不同组件层级的系统开销。

\begin{table}[htbp]
    \centering
    \small
    \caption{基准工作负载的执行特征与主导成本分析}
    \label{tab:workload-cost-types}
    \begin{tabular}{p{2.0cm}p{3.2cm}p{5.1cm}p{1.9cm}}
        \hline
        \textbf{工作负载} & \textbf{类型与生命周期} & \textbf{主导成本观察重点} & \textbf{相关章节} \\
        \hline
        \texttt{pmd} & DaCapo，短生命周期、控制流密集 & 解释执行驻留、JIT 触发时机、编译固定开销 & 5.2节 \\
        \texttt{mpegaudio} & SPEC 吞吐型、运行时热路径对照 & 运行时数据传输、字符串及库函数调用边界 & 5.3节 \\
        \texttt{lu.small} & SPEC 规则浮点数值内核 & 向量化命中率、LSX 映射收益与启用边界 & 第 6 章 \\
        \texttt{helloworld} & SPEC 启动型、加载期敏感 & DEX 校验及容器探测开销、机制解耦收益 & 5.4节 \\
        \hline
    \end{tabular}
\end{table}

表~\ref{tab:workload-cost-types} 表明，本文探讨的优化案例重点对应三类核心执行开销：短生命周期应用的控制流开销、高频运行时的访存数据流开销，以及规则数值计算的数据并行开销。此外，\texttt{sor.small} 与 \texttt{fft.small} 将在后文被用作性能对比验证集，以评估向量化优化方案的泛化能力。

\section{本章小结}

本章围绕执行链路构建与性能瓶颈分析，梳理了平台适配的相关技术基础及架构约束。首先，ART 虚拟机的跨架构移植需构建涵盖解释执行、动态编译与运行时辅助机制的完整执行生态。其次，LoongArch 硬件特征中的控制流组织规则、底层内存访存代价及硬件级向量扩展，共同构成了制约动态执行性能的关键瓶颈维度。最后，优化目标的定位高度依赖于具体工作负载的生命周期与主导开销类型。后续章节将依托上述理论框架，开展执行系统底座的指令级实现，并在统一实验环境中对相应的优化策略展开定量验证与边界评估。